\section{Results I: Parameter Count}
\label{sec:results-params}

Table~\ref{tab:configs} (Section~\ref{sec:configs}) already lists each configuration's classical and
quantum parameter count; Table~\ref{tab:paramshare} restates the same counts as the two quantities that
matter for the parameter-count claim -- the quantum branch's share of the total, and the reduction relative
to the classical baseline.

\begin{table}[h]
\centering
\caption{Quantum share of the total parameter count, and the reduction relative to the classical baseline
(4353 parameters). Quantum share is $4nL$ divided by total parameters.}
\label{tab:paramshare}
\begin{tabular}{@{}lrr@{}}
\toprule
Config & Quantum share & Reduction vs.\ classical \\
\midrule
narrow   & 0.66\% & 44.4\% \\
baseline & 1.25\% & 41.0\% \\
wide     & 1.77\% & 37.6\% \\
deep     & 2.46\% & 40.2\% \\
\bottomrule
\end{tabular}
\end{table}

Across the four hybrid configurations, the VQC itself contributes only 16--64 trainable parameters -- at
most 2.46\% of the network's total. Total parameter count nonetheless falls by 37.6--44.4\% relative to the
classical baseline, because the classical pre- and post-processor around the VQC is already narrower than
the classical baseline's five hidden layers of width 32 (Section~\ref{sec:models}). The headline
parameter-count advantage of the hybrid models is therefore almost entirely a property of the classical
architecture chosen around the quantum layer, not of the quantum layer itself. This is what motivates the
size-matched controls of Section~\ref{sec:setup-controls}: comparing a hybrid only to the classical baseline
would credit the VQC with parameter savings it did not produce.
